\documentclass[12pt, a4paper]{article}

\usepackage[utf8]{inputenc}
\usepackage[T2A]{fontenc}
\usepackage[russian]{babel}
\usepackage{amsmath, amssymb, amsfonts}
\usepackage{booktabs}
\usepackage{tabularx}
\usepackage{array}
\usepackage{geometry}
\geometry{left=30mm, right=15mm, top=20mm, bottom=20mm}
\usepackage{setspace}
\onehalfspacing
\setlength{\emergencystretch}{3em}
\usepackage{caption}
\usepackage{float}
\usepackage{graphicx}
\usepackage{indentfirst}
\usepackage{enumitem}
\usepackage{cite}

\graphicspath{{assets/}}

\captionsetup{font=small, labelsep=period}
\captionsetup[figure]{name=Рисунок}

\setlist[itemize]{itemsep=0pt, topsep=0pt, parsep=0pt, partopsep=0pt}
\setlist[enumerate]{itemsep=0pt, topsep=0pt, parsep=0pt, partopsep=0pt}

% Рисунок на всю строку с подписью и источником
\newcommand{\figurereport}[4]{
\begin{figure}[H]
    \centering
    \includegraphics[width=#2\textwidth]{#1}
    \caption{#3}
    \vspace{-2mm}
    {\footnotesize Источник: #4.}
\end{figure}
}

\newcommand{\tablesource}[1]{\par\vspace{1mm}{\footnotesize Источник: #1.}}

\begin{document}

\begin{titlepage}
    \thispagestyle{empty}
    \begin{center}
        \linespread{1.1}\selectfont
        \textbf{Домашнее задание №1}\\[1cm]
        \Large{\textbf{Пространственная эконометрика}}\\[1cm]
        \large{Пространственные признаки в прогнозировании\\
        розничного товарооборота магазинов «Пятерочка»}\\[0.6cm]
        \normalsize{Воспроизведение исходного ансамбля и результаты эксперимента}\\[1.5cm]
        \normalsize
        \textbf{Выполнили:}\\[0.4cm]
        Гончаров Алексей\\
        Дугарев Денис\\
        Русских Никита\\
        Смирнов Владимир\\
        Хмельницкий Богдан\\[2cm]
    \end{center}
    \vfill
    \begin{center}
        Москва 2026
    \end{center}
\end{titlepage}

\setcounter{page}{2}

% ============================================================
\section*{Постановка задачи и данные}

Цель эксперимента - проверить, могут ли пространственные признаки повысить точность прогноза розничного товарооборота (РТО) магазинов «Пятерочка».
В основной части исследования пространственные связи изучались эконометрическими методами.
Здесь использован более прикладной подход; характеристики соседних населенных пунктов добавлены к признакам моделей градиентного бустинга,
после чего сравнивается качество прогнозов.

В исходных данных 485\,082 наблюдения за январь 2023 - февраль 2025 года.
Каждая строка соответствует магазину в определенном месяце; задача - спрогнозировать РТО за март 2025 года.
Соревновательная оценка основана на средней абсолютной процентной ошибке (MAPE):
\begin{equation}
\operatorname{MAPE}=\frac{100}{n}\sum_{i=1}^n\left|\frac{\widehat y_i-y_i}{y_i}\right|,
\qquad
S=100\left(1-\frac{\min(\operatorname{MAPE},100)}{100}\right)^2.
\label{eq:score}
\end{equation}

В качестве отправной точки взят исходный ансамбль четырех моделей градиентного бустинга над решающими деревьями \cite{friedman2001}.
Он получил 90,64 балла на публичной части теста и 90,78 балла на полном итоговом тесте.
Все результаты нового эксперимента далее приводятся именно для полного теста, если не указано иное. В таблице
ниже отдельно указаны результаты для обогащенного пространственными признаками ансамбля из градиентных бустингов и несколько иным видом
агрегации прогнозов базовых моделей.

\begin{table}[H]
\centering
\footnotesize
\caption{Сравнение исходного и итогового результатов}
\begin{tabularx}{\textwidth}{Xrr}\toprule
Вариант & Публичная часть & Полный тест\\\midrule
Исходный ансамбль & 90,64 & 90,78\\
Ансамбль с пространственными признаками & - & \textbf{90,83}\\\bottomrule
\end{tabularx}
\tablesource{расчеты авторов по данным X5 Group}
\end{table}

\subsection*{Исходное решение: отбор признаков и ансамблирование}

Процесс отбора признаков здесь достаточно непримечателен, хотя, в двух словах отметить стоит.
Для исходного признакового набора генерируется порядка 380 фичей (скользящие окна, лаги, агрегаты и прочие).
Далее проводится своего рода скорректированный под особенности задачи Forward Selection с нуля - считаем важность по
лоссу базовой модели; например, CatBoost глубиной 5 в качестве базовой модели для всего признакового пространства,
а после итеративно отбираем по одной фиче из отсортированного по важности по вкладу в лосс в пул выбранных фичей.
Каждые несколько фичей важность исходного набора (уже из оставшихся признаков) пересчитывается и процесс отбора сверху вниз запускается заново
(это что-то вроде страховки против дублирования сигналов, а также возможность учета перекрестных значимостей).
Само значение вклада в лосс считается усреднением между 5 различными фолдами-месяцами в данных с различными весами.

Процесс подбора гиперпараметров аналогично проводился по взвешенному значению метрики на валидации,
но уже на 3-х фолдах (использовалась Optuna, метод - multivariate TPE) \cite{akiba2019, bergstra2011}.

Но более интересна здесь все же сама модель.

Подсвечу, что отбор признаков происходил для разных моделей вида градиентного бустинга;
под разными имею в виду обучение на разных таргетах - логарифм розничного товарооборота и лог-дифференцированное значение розничного товарооборота.
Соответственно, признаки были отобраны несколько разные для каждого из бустингов, но с некоторыми пересечениями.

После получения прогнозов на тестовых данных значения прогнозов усреднялись 4-мя различными методами усреднения, а после еще раз усреднялись
простым арифметическим средним для итогового предикта.

В общем, фактически, был реализован частный случай стекинга - блендинг с некоторой добавленной стохастичностью. Почему это работает?

Вспомним, что чем менее скоррелированы между собой ошибки базовых алгоритмов, тем сильнее мы снизим дисперсию итогового прогноза:

\[
E_N = E_1\left(\rho + \frac{1-\rho}{N}\right).
\]

Или же словами:

\textbf{Дисперсия ошибки композиции (ансамбля) из N моделей = Дисперсия ошибки одной из отдельных моделей *
(коэффициент корреляции + (1 - коэф. корреляции) / N), где коэффициент корреляции ответственен за корреляцию между
ошибками базовых моделей}

В задаче, справедливости ради, прогнозы были достаточно схожими
(то есть, высокая корреляция перебивала полезный эффект от N моделей), поэтому эффект не был в достаточной мере большим,
но так или иначе среднее сглаживало перепрогнозы и недопрогнозы отдельных моделей.

Сама формула, скорее, схематично показывает причину улучшения итогового скора на лидерборде, так как,
с большой вероятностью, при ансамблировании ансамблей она становится заметно сложнее, но, тем не менее, по сути является таковой в обычном бэггинге.

Может ли быть оптимальное число моделей в таком ансамбле? Обратимся к комбинаторной теории переобучения:

\begin{itemize}\item Если у нас есть хорошая модель, а остальные - плохие, то хорошая модель будет превалировать и мы не будем переобучаться;

\item Если есть много похожих друг на друга хороших моделей, мы тоже не будем переобучаться, так как эффективная сложность
совокупности похожих друг на друга моделей невелика;

\item Но если получится так, что модели существенно различные и
все примерно одинаково плохие, то переобучение будет очень сильным, а сам эффект переобучения растет по мере роста числа моделей \cite{vorontsov2016}.
\end{itemize}

Отсюда заметим, что, в целом, если все модели достаточно хорошие, имеет смысл увеличивать их количество; тем не менее,
на практике достаточно понять на валидации оптимальное минимальное число.

Прикрепляю схему алгоритма. Так немного проще визуально понять, что именно и каким образом реализуется.

\figurereport{model_scheme.jpg}{0.88}{Схема алгоритма исходного ансамбля}{составлено авторами}

\subsection*{Состав моделей и правило объединения прогнозов}

В новом эксперименте исходные признаки, преобразования целевой переменной и параметры обучения сохранены.
Это позволяет отдельно оценить эффект пространственных признаков, не смешивая его с результатом дополнительного подбора гиперпараметров.

\begin{table}[H]
\centering
\footnotesize
\caption{Модели исходного ансамбля}
\begin{tabularx}{\textwidth}{Xlr}\toprule
Модель & Целевая переменная & Число деревьев\\\midrule
CatBoost \cite{prokhorenkova2018} & $\log(1+y)$ & 6500\\
LightGBM (приросты) \cite{ke2017} & $\Delta\log y$ & 7000\\
LightGBM (уровни) & $\log(1+y)$ & 8000\\
XGBoost (приросты) \cite{chen2016} & $\Delta\log y$ & 3750\\\bottomrule
\end{tabularx}
\tablesource{составлено авторами}
\end{table}

Прогнозы моделей, обученных на логарифме РТО, возвращаются к исходной шкале с помощью $\exp(\widehat z)-1$.
Для моделей, прогнозирующих логарифмический прирост, применяется $y_{i,t-1}\exp(\widehat g_{it})$, где $g_{it}=\log y_{it}-\log y_{i,t-1}$.
Ошибка обучения рассчитывается для преобразованной целевой переменной, а итоговая MAPE - для РТО в исходной шкале.

Для четырех прогнозов $p_1,\dots,p_4$ арифметическое и логарифмическое средние имеют вид:
\begin{equation}
A=\frac14\sum_{m=1}^4 p_m, \qquad
L=\exp\left(\frac14\sum_{m=1}^4\log(1+p_m)\right)-1.
\end{equation}

Дополнительно используются два пороговых правила, обозначенные $G$ и $H$. В $G$
арифметическое среднее заменяется медианой, когда $\max p_m/\min p_m>1{,}08$; в $H$ действует аналогичное правило,
но без модели LightGBM, прогнозирующей приросты. В принятой реализации итоговый прогноз равен:
\begin{equation}
P=\exp\left(\frac14\sum_{q\in\{A,L,G,H\}}\log(1+q)\right)-1.
\label{eq:meta}
\end{equation}
Именно это правило объединения прогнозов было сохранено при основном сравнении составов ансамбля.

\clearpage

% ============================================================
\section*{Построение пространственных признаков}

Основная идея взята из пространственной модели SLX \cite{vega2015}, в которой наряду с собственными характеристиками объекта используются характеристики его соседей:
\begin{equation}
y=X\beta+WX\theta+\varepsilon.
\end{equation}
Здесь $W$ - матрица пространственных весов \cite{lesage2009}, а $WX$ - взвешенные характеристики соседних территорий. В отличие от линейной модели SLX,
в нашем эксперименте $WX$ используется только как дополнительный набор входных признаков бустинга.
Прямые и косвенные причинные эффекты не оцениваются.

Прогноз строился для 18\,657 магазинов в 3\,360 населенных пунктах.
Пространственные признаки рассчитывались на уровне населенного пункта, а затем присоединялись к наблюдениям отдельных магазинов.
Таким образом, исходная панель магазинов не заменялась агрегированными данными.

Для построения $W$ проверены два варианта соседства: по 6 и по 12 ближайшим населенным пунктам (kNN6 и kNN12).
Расстояния между пунктами рассчитывались в метрической проекции Альберса с эллипсоидом Красовского.
Если один населенный пункт входит в список ближайших соседей другого, связь учитывается в обе стороны.
Затем веса нормируются так, чтобы их сумма для каждой строки равнялась единице:
\begin{equation}
w_{rv}=\frac{a_{rv}}{\sum_u a_{ru}}, \qquad w_{rr}=0,\qquad \sum_v w_{rv}=1.
\end{equation}

Для числового признака сначала находится среднее по магазинам населенного пункта, затем - взвешенное среднее по соседним пунктам:
\begin{equation}
\overline x_{rjt}=\frac{1}{|\mathcal I_{rj}^{\mathrm{obs}}(t)|}
\sum_{i\in\mathcal I_{rj}^{\mathrm{obs}}(t)}x_{ijt}, \qquad
(WX)_{rjt}=\sum_v w_{rv}\overline x_{vjt}.
\end{equation}
Магазины одного населенного пункта не считаются соседями друг для друга при расчете $WX$.
Для каждой модели формируется свой набор пространственных признаков на основе доступных исходных переменных;
идентификаторы и технические поля в него не включаются.
В итоге к CatBoost было добавлено 39 признаков, к LightGBM приростов - 22, к LightGBM уровней - 33, к XGBoost - 14.
Часть общих календарных признаков после пространственного усреднения остается неизменной и не дает нового сигнала.

Для прогноза марта использовались только сведения, доступные к концу февраля.
Фактический мартовский РТО не применялся при построении признаков; исторические лаги, напротив, были допустимы.
Это важно, чтобы избежать утечки целевой переменной. Пропущенные значения в пространственных признаках сохранялись без дополнительного заполнения.

\clearpage

% ============================================================
\section*{Эксперименты и результаты}

\subsection*{Выбор пространственных признаков на исторических данных}

Для каждой из четырех моделей были проверены исходная версия без $WX$ и два варианта с kNN6 и kNN12.
Сначала варианты сравнивались по средней MAPE за август, ноябрь и декабрь 2024 года.
Затем выбранные пространственные версии и исходные модели проверялись на январе и феврале 2025 года.
В каждом случае обучение проводилось только на предшествующих месяцах.

\begin{table}[H]
\centering
\footnotesize
\caption{MAPE на этапе отбора пространственных признаков, \%}
\begin{tabularx}{\textwidth}{Xrrrr}\toprule
Модель & Без $WX$ & kNN6 & kNN12 & Выбор\\\midrule
CatBoost, уровни & 4,333 & 4,238 & 4,217 & $k=12$ \\
LGB, приросты & 5,335 & 4,179 & 4,080 & $k=12$ \\
LGB, уровни & 3,766 & 3,786 & 3,812 & $k=6$ \\
XGBoost, приросты & 4,864 & 4,744 & 4,664 & $k=12$ \\
\bottomrule\end{tabularx}
\tablesource{расчеты авторов по данным X5 Group и GeoNames}
\end{table}

\begin{table}[H]
\centering
\footnotesize
\setlength{\tabcolsep}{4pt}
\caption{MAPE на проверке в январе и феврале 2025 года, \%}
\begin{tabularx}{\textwidth}{Xrrrrrr}\toprule
& & & \multicolumn{2}{c}{Январь} & \multicolumn{2}{c}{Февраль}\\
\cmidrule(lr){4-5}\cmidrule(lr){6-7}
Модель & Деревья & $k$ & До & С $WX$ & До & С $WX$\\\midrule
CatBoost, уровни & 6\,500 & 12 & 4,276 & 3,961 & 4,264 & 4,448 \\
LGB, приросты & 7\,000 & 12 & 4,003 & 4,000 & 5,122 & 4,689 \\
LGB, уровни & 8\,000 & 6 & 4,179 & 4,242 & 4,645 & 5,169 \\
XGBoost, приросты & 3\,750 & 12 & 4,030 & 4,064 & 4,428 & 4,714 \\
\bottomrule\end{tabularx}
\tablesource{расчеты авторов по данным X5 Group и GeoNames}
\end{table}

Здесь результаты получились неоднозначными. CatBoost стал точнее в январе, но хуже в феврале.
LightGBM, прогнозирующий логарифмические приросты, улучшил MAPE в обоих месяцах; в феврале ошибка снизилась с 5,122 до 4,689\%.
Два других бустинга - LightGBM уровней и XGBoost - на этих месяцах уступили исходным версиям.
Поэтому по предварительному критерию для замены в ансамбле подходил только LightGBM приростов.

\subsection*{Результаты отдельных моделей на мартовском тесте}

\begin{table}[H]
\centering
\footnotesize
\caption{Итоговый балл отдельных моделей до и после добавления $WX$}
\begin{tabularx}{\textwidth}{Xrrr}\toprule
Модель & Без $WX$ & С $WX$ & Изменение\\\midrule
CatBoost, уровни & 90,55 & 90,39 & -0,16 \\
LGB, приросты & 90,32 & 90,24 & -0,08 \\
LGB, уровни & 90,57 & 90,55 & -0,02 \\
XGBoost, приросты & 90,45 & 90,58 & 0,13 \\
\bottomrule\end{tabularx}
\tablesource{расчеты авторов по данным X5 Group и GeoNames}
\end{table}

На полном тесте марта результат улучшился только у XGBoost: с 90,45 до 90,58 балла.
CatBoost потерял 0,16 балла, LightGBM приростов - 0,08, LightGBM уровней - 0,02.
Следовательно, преимущество LightGBM приростов на январе и феврале не перенеслось на март, а XGBoost, наоборот, показал улучшение.

\figurereport{model_comparison.png}{0.92}{Изменение ошибок на исторических месяцах и баллов на мартовском тесте}{составлено авторами по данным X5 Group и GeoNames}

Эти результаты наглядно показывают, что качество отдельных моделей на исторических месяцах не всегда соответствует их
результату на итоговом тестовом наборе данных в рамках соревнования. При этом важно оценивать не только отдельные прогнозы, но и их сочетание в ансамбле.

\subsection*{Сравнение составов ансамбля}

Чтобы сравнение было корректным, формула объединения прогнозов (\ref{eq:meta}) не менялась.
В каждом варианте заменялись только отдельные исходные модели на версии с пространственными признаками.

\begin{table}[H]
\centering
\footnotesize
\setlength{\tabcolsep}{4pt}
\caption{Сравнение составов ансамбля}
\begin{tabularx}{\textwidth}{Xrrr}\toprule
Модели с пространственными признаками & Январь, \% & Февраль, \% & Балл за март\\\midrule
Нет (исходный ансамбль) & 3,811945 & 4,265305 & 90,78\\
Только LightGBM приростов & 3,807809 & 4,209390 & 90,78\\
Все четыре модели & 3,763799 & 4,263002 & 90,79\\
Только XGBoost & 3,796333 & 4,291810 & 90,80\\
CatBoost и XGBoost & 3,754268 & 4,209174 & 90,76\\
XGBoost и LightGBM приростов & 3,789656 & 4,240780 & 90,82\\
XGBoost и LightGBM уровней & 3,805225 & 4,402870 & 90,81\\
XGBoost и оба LightGBM & 3,799520 & 4,349020 & \textbf{90,83}\\\bottomrule
\end{tabularx}
\tablesource{расчеты авторов по данным X5 Group и GeoNames}
\end{table}

Предварительно выбранная замена LightGBM приростов уменьшила среднюю MAPE января и февраля с 4,038625 до 4,008600\%,
но на мартовском тесте балл остался равным 90,78. Сочетание CatBoost и XGBoost с $WX$, лучшее среди приведенных вариантов
по историческим месяцам, получило только 90,76.

Лучший мартовский результат, 90,83, получен при сохранении исходного CatBoost и замене XGBoost и двух LightGBM на версии с $WX$.
При этом средняя MAPE января и февраля у этого ансамбля составила 4,074270\% против 4,038625\% у исходного.
Таким образом, отбор по январю и февралю не привел бы к выбору именно этого состава.

Важно, что замена LightGBM, взятая отдельно, могла ухудшить его балл, но оказаться полезной в ансамбле.
По сути, это не является противоречием; вклад модели в итоговый прогноз зависит также от ошибок других компонентов.

\subsection*{Дополнительное сравнение способов усреднения}

\figurereport{ensemble_grid.png}{0.92}{Результаты вариантов усреднения на мартовском тесте}{составлено авторами по данным X5 Group и GeoNames}

Максимальные 90,83 балла получили два варианта:
исходное правило ансамблирования с $WX$ у XGBoost и обоих LightGBM (далее $A$) и вариант той же четверки
с пороговым выбором медианы при расхождении прогнозов более 1,08 (далее $B$). Простое арифметическое и логарифмическое
средние этих четырех моделей дали по 90,81 балла.

Также были рассмотрены два состава с результатом 90,82: $C$ - с $WX$ у XGBoost и
LightGBM приростов и пороговым правилом; $D$ - с $WX$ у XGBoost и LightGBM уровней, без LightGBM приростов.
Заключительная проверка касалась пяти способов объединения готовых ансамблей:

\begin{table}[H]
\centering
\footnotesize
\caption{Сравнение комбинаций готовых ансамблей}
\begin{tabularx}{\textwidth}{rXr}\toprule
№ & Способ объединения & Балл за март\\\midrule
1 & Арифметическое среднее $A,B$ & 90,83\\
2 & Логарифмическое среднее $A,B$ & 90,83\\
3 & Арифметическое среднее $A,B,C$ & 90,83\\
4 & Арифметическое среднее $A,B,D$ & 90,83\\
5 & Логарифмическое среднее $A,B,D$ & 90,83\\\bottomrule
\end{tabularx}
\tablesource{расчеты авторов по данным X5 Group и GeoNames}
\end{table}

Ни один вариант не превысил результат 90,83 балла после округления. Для окончательного решения оставлен ансамбль $A$:
он обеспечивает тот же результат без дополнительного уровня усреднения.
Однако по опубликованным округленным оценкам нельзя утверждать, что его точная ошибка минимальна среди всех проверенных вариантов.

\clearpage

% ============================================================
\section*{Пространственная зависимость ошибок и ограничения}

Для февраля 2025 года дополнительно рассчитан индекс Морана \cite{moran1950} по средним относительным ошибкам магазинов внутри населенных пунктов:
\begin{equation}
e_r=\operatorname{mean}_{i\in r}\left[\frac{y_i-\widehat y_i}{y_i}\right].
\end{equation}
Значимость проверялась перестановочным тестом с 999 перестановками. Здесь, в отличие от MAPE по магазинам, каждая территория имеет одинаковый вес.

\begin{table}[H]
\centering
\footnotesize
\caption{Индекс Морана для ошибок прогноза, февраль 2025 года}
\begin{tabularx}{\textwidth}{Xrrrr}\toprule
Модель & $k$ & До & С $WX$ & $p$-значение\\\midrule
CatBoost, уровни & 12 & 0,0657 & 0,0565 & 0,001 \\
LGB, приросты & 12 & 0,0559 & 0,0524 & 0,001 \\
LGB, уровни & 6 & 0,0960 & 0,0980 & 0,001 \\
XGBoost, приросты & 12 & 0,0710 & 0,0470 & 0,001 \\
\bottomrule\end{tabularx}
\tablesource{расчеты авторов по данным X5 Group и GeoNames}
\end{table}

Пространственная зависимость ошибок сохраняется. Например, у XGBoost индекс Морана уменьшился с 0,0710 до 0,0470,
хотя февральская MAPE выросла с 4,428 до 4,714\%. Следовательно, уменьшение пространственной автокорреляции само по себе
еще не означает улучшения точности прогноза.

До основного эксперимента также рассматривались признаки на основе истории соседей, локальных пространственных статистик
и собственных векторов матрицы весов, а также варианты пространственных эконометрических моделей.
Некоторые из них улучшали показатели на отдельных исторических периодах, но не дали устойчивого выигрыша на мартовском тесте.
В отдельной проверке ранняя остановка при обучении также ухудшила качество относительно исходного режима.
Поэтому эти варианты не вошли в окончательный ансамбль.

Есть важное ограничение: после получения результатов полного мартовского теста они использовались при сравнении составов ансамбля.
Следовательно, 90,83 балла - результат на конкретном уже проверенном тесте, а не независимая оценка качества на будущих месяцах.
Кроме того, без истинных значений мартовского РТО нельзя изучить корреляции ошибок базовых моделей и проверить статистическую
устойчивость полученного улучшения.

\clearpage

% ============================================================
\section*{Выводы}

Добавление пространственных признаков позволило повысить результат исходного ансамбля с \textbf{90,78 до 90,83 балла}
на полном мартовском тесте. По формуле соревновательной метрики это приблизительно соответствует снижению MAPE с 4,7215 до 4,6952\%,
то есть на \textbf{0,0262 процентного пункта}, или примерно на 0,56\% относительно исходной ошибки.
Значения MAPE приближенные, поскольку баллы опубликованы с округлением.

При улучшении итогового результата с \textbf{90,78 до 90,83 балла} команда поднялась бы с \textbf{21-22-го на 12-е место}
в общем рейтинге и с \textbf{8-го на 5-е место} среди команд, приглашенных на защиту.
Таким образом, даже относительно небольшое улучшение качества прогноза позволило бы существенно повысить позицию команды в
итоговом рейтинге.

Наиболее полезной отдельной пространственной версией оказалась XGBoost. Но в лучшем ансамбле были также заменены
оба LightGBM, хотя по отдельности их баллы ухудшились. CatBoost остался без пространственных добавок. Эксперименты
с дополнительным усреднением не улучшили округленный итоговый результат.

Таким образом, пространственные характеристики соседних населенных пунктов могут давать дополнительный сигнал для прогнозирования РТО,
но эффект в этом эксперименте оказался небольшим и зависел от состава ансамбля; тем не менее, как таковой выигрыш от новых признаков заметен.

\subsection*{Материалы эксперимента}

Для воспроизведения результатов использованы исходные данные и ноутбук конкурсного решения, сохраненные прогнозы и
материалы семинаров по пространственной эконометрике. Рабочие файлы, таблицы и результаты проверок собраны в
отдельном каталоге проекта. Исходные файлы основного эконометрического исследования не изменялись.

\clearpage

\begin{thebibliography}{99}
    \bibitem{vorontsov2016}
    Воронцов, К. В. Математические методы прогнозирования объемов продаж / К. В. Воронцов // ПостНаука : канал YouTube. - 2016. - 25 февраля. - URL: https://www.youtube.com/watch?v=u3lhVV3wAKk (дата обращения: 13.10.2026).

    \bibitem{chen2016}
    Chen, T. XGBoost: A Scalable Tree Boosting System / T. Chen, C. Guestrin // Proceedings of the 22nd ACM SIGKDD International Conference on Knowledge Discovery and Data Mining. - 2016. - P. 785-794.

    \bibitem{friedman2001}
    Friedman, J. H. Greedy Function Approximation: A Gradient Boosting Machine / J. H. Friedman // Annals of Statistics. - 2001. - Vol. 29, No. 5. - P. 1189-1232.

    \bibitem{vega2015}
    Halleck Vega, S. The SLX Model / S. Halleck Vega, J. P. Elhorst // Journal of Regional Science. - 2015. - Vol. 55, No. 3. - P. 339-363.

    \bibitem{ke2017}
    LightGBM: A Highly Efficient Gradient Boosting Decision Tree / G. Ke, Q. Meng, T. Finley, T. Wang, W. Chen, W. Ma, Q. Ye, T.-Y. Liu // Advances in Neural Information Processing Systems. - 2017. - Vol. 30. - P. 3146-3154.

    \bibitem{lesage2009}
    LeSage, J. Introduction to Spatial Econometrics / J. LeSage, R. K. Pace. - Boca Raton : Chapman and Hall/CRC, 2009.

    \bibitem{moran1950}
    Moran, P. A. P. Notes on Continuous Stochastic Phenomena / P. A. P. Moran // Biometrika. - 1950. - Vol. 37, No. 1/2. - P. 17-23.

    \bibitem{akiba2019}
    Optuna: A Next-generation Hyperparameter Optimization Framework / T. Akiba, S. Sano, T. Yanase, T. Ohta, M. Koyama // Proceedings of the 25th ACM SIGKDD International Conference on Knowledge Discovery and Data Mining. - 2019. - P. 2623-2631.

    \bibitem{bergstra2011}
    Algorithms for Hyper-Parameter Optimization / J. Bergstra, R. Bardenet, Y. Bengio, B. Kegl // Advances in Neural Information Processing Systems. - 2011. - Vol. 24. - P. 2546-2554.

    \bibitem{prokhorenkova2018}
    CatBoost: Unbiased Boosting with Categorical Features / L. Prokhorenkova, G. Gusev, A. Vorobev, A. V. Dorogush, A. Gulin // Advances in Neural Information Processing Systems. - 2018. - Vol. 31. - P. 6638-6648.
\end{thebibliography}
\end{document}
